\documentclass[12pt]{article}
\usepackage[utf8]{inputenc}
\usepackage[T1]{fontenc}
\usepackage[margin=1in]{geometry}
\usepackage{setspace}
\usepackage{xcolor}
\usepackage[colorlinks=true, urlcolor=blue!60!black, linkcolor=blue!60!black]{hyperref}
\setlength{\parindent}{0em}
\setlength{\parskip}{0.6em}
\renewcommand{\baselinestretch}{1.0}
\begin{document}
\begin{center}
{\LARGE \textit{A.L.A. Schechter Poultry Corp. v. United States}, 295 U.S. 495 (1935)\par}
\vspace{0.5em}
{\large Case Brief --- Two-Page Edition\par}
\end{center}

\section{Memory Jogger}
In the ``sick chicken case,'' a unanimous Court held that NIRA \S~3 unconstitutionally delegated lawmaking power to the President and trade groups and that regulating a Brooklyn slaughterhouse's local business exceeded the commerce power (\href{https://supreme.justia.com/cases/federal/us/295/495/}{295 U.S. 495}).

\section{Detailed Case Facts}
The National Industrial Recovery Act let the President approve ``codes of fair competition'' proposed by trade associations, enforceable as crimes. The Live Poultry Code for New York set wages and hours and a ``straight killing'' rule. The Schechter brothers bought poultry shipped from other States, slaughtered it in Brooklyn, and sold it only to local retailers.

\section{Procedural History}
Convicted on 19 counts, including selling an unfit chicken, the Schechters won reversal in the Second Circuit only on the wage-and-hour counts. Both sides sought review.

\section{Judicial Votes}
9--0. Hughes wrote for the Court; Cardozo, joined by Stone, concurred.

\section{Holding}
Section 3 unconstitutionally delegates legislative power by authorizing binding codes without standards, and the code's application to local slaughtering and sales exceeds the commerce power.

\section{Analysis of Opinions}
Hughes held that ``[e]xtraordinary conditions do not create or enlarge constitutional power.'' Congress must set policy and standards, but NIRA never defined ``fair competition,'' handed proposals to private groups in a way ``unknown to our law,'' and left the President virtually unfettered. The poultry had come to ``a permanent rest within the State,'' and its local handling affected commerce only indirectly. Cardozo, who had dissented in \textit{Panama Refining}, called \S~3 a ``roving commission'' and warned that without limits of degree ``there will be an end to our federal system.''

\section{Examples: Future Applications}
\textbf{Same-Side:} criminally enforced AI codes drafted by leading labs, without APA procedures, are invalid. \textbf{Opposite-Side:} health-based air standards set by an agency through rulemaking supply an intelligible principle (\textit{Whitman}). \textbf{Fence-Sitter:} criminally enforced agency rules to promote ``fair and orderly markets.''

\section{Critique}
\textit{Schechter} reflects suspicion of corporatist self-regulation that let dominant firms cartelize industries. Justice Scalia noted that ``how much is too much'' is a question of degree ill suited to courts, and the absence of procedural safeguards made NIRA uniquely dangerous.

\section{Key Quotations}
\begin{quote}
``The delegated power of legislation which has found expression in this code is not canalized within banks that keep it from overflowing. It is unconfined and vagrant \dots\ This is delegation running riot. No such plenitude of power is susceptible of transfer.''
\end{quote}
\noindent --- Justice Cardozo, concurring, 295 U.S. at 551, 553.

\end{document}
